\section{Relevanz}

Kleine und mittlere Unternehmen bilden das Rückgrat der europäischen
Unternehmenslandschaft. Rund 99\,\% aller Unternehmen in der Europäischen
Union zählen zu den KMU \parencite{Alaskari2021,Leyh2015}. Damit ist die
digitale Transformation von KMU nicht nur eine einzelbetriebliche
Herausforderung, sondern auch von gesamtwirtschaftlicher Bedeutung.
Gleichzeitig lassen sich Ansätze und Erfahrungen aus Großunternehmen nicht
ohne Weiteres auf KMU übertragen
\parencite{Xia2009,Leyh2015}. 

Der Handlungsdruck wird durch einen stetigen Wandel der
Verbraucheranforderungen zusätzlich verstärkt \parencite{Alaskari2021}.
Kunden erwarten zunehmend schnelle Reaktionszeiten, transparente Informationen
und flexible Leistungen. 
Für KMU entsteht dadurch die Notwendigkeit ihre Prozesse effizienter zu gestalten \parencite{Leyh2017}.

ERP-Systeme können hierfür eine zentrale technologische Grundlage bilden \parencite{Leyh2017}. 
Deren Nutzen entsteht jedoch
erst dann, wenn das System wirksam in die organisatorischen,
technologischen und prozessualen Bedingungen eines Unternehmens eingebettet
wird.

Die Relevanz des Themas ergibt sich somit aus der hohen wirtschaftlichen
Bedeutung von KMU, dem steigenden digitalen Anpassungsdruck und der zentralen
Rolle von ERP-Systemen als möglicher Grundlage integrierter, datenbasierter
Unternehmensprozesse.